\section{Definite theta periods and the exact iterated seesaw}

There are now two precise global comparisons. First, an orthogonal decomposition of the rank-three Hermitian space gives a totally definite rank-two subgroup. Restriction to that subgroup is a genuine global unitary Bessel period, and the theta tensor rule gives an exact seesaw identity; this is proved in THETA3-B11--B12. Second, the recovered Alp\"oge source supplies a common normalized functional for an entire family of iterated lifts. Its exact statement is \texttt{MainX}, and its source witness and elementary scalar consequence are recorded in THETA3-B13. Identifying the particular cohomological inputs with that family, including its finite lattices and normalization, remains an instantiation problem.

The inner-form calculations below describe a different possible route. Their obstructions do not obstruct the definite-subgroup construction. All new proofs have been read in full for the targeted checks recorded in the accompanying event files. Source statements used as background are distinguished from these proofs. No Lean compilation, complete dependency audit, or validation of the cited RH theorem is claimed or attempted.

\subsection{THETA3-B0: what the supplied and recovered sources establish}\label{B0}

The attached four-page Alp\"oge note gives a timeline. Step 4 mentions determinant positivity, step 7 GGP periods and interlacing, step 8 iterated theta lifts on totally definite unitary groups, and step 12 a seesaw producing a Bessel period independent of the Hecke characters. The note does not specify the integral, test vectors, measures, or comparison with differential forms. The attached source map is explicitly a text-parsing report, with no Lean build. Its paths are useful search locators; they are not substitutes for the definitions at those paths. The attached \texttt{Main.lean} and \texttt{Statements.lean} contain the final analytic declarations, not the theta-period definitions. This is an evidence finding, not a refutation of the intended comparison.

During this run, the named definitions and several proof bodies were recovered from the secondary public mirror \texttt{CocoSgt/\allowbreak{}riemann-hypothesis-lean}, pinned to commit
\begin{center}\small\texttt{7cbf7a0c32fca1badf05931879a37c3f8a2b03b4}.\end{center}
The commit date is 20 November 2026. The recovered \texttt{Main.lean} and \texttt{Statements.lean} are byte-for-byte equal to the attachments. A manifest records the pinned URL, SHA256, Git blob hash, size, and retrieval time of each fetched file. This is partial identity evidence; it does not assert a complete primary-archive hash match. Consequently the initial missing-definition gate has been reduced to the exact mathematical matching problem in B13.

For reference, the 032 setting is a Galois CM extension $E/F$, with $[F:\Q]\geq2$, and an $E$-Hermitian space $V$ of dimension three. At one real place $\mu$ it has signature $(2,1)$, and elsewhere signature $(3,0)$. The compact ball surface is attached to $\SU(V)$ at sufficiently small principal level. For multiplicity space $P$ of dimension $k=1$ or $2$, the alternating form is
\[
 \psi_P=\Tr_{E/F}\bigl(\delta h_V\otimes h_P\bigr),
 \qquad \overline\delta=-\delta.
\]
At the active place the positive Lagrangian is
\[
 \mathcal F_P=(L\otimes P_\mu)\oplus(Q^*\otimes P_\mu^*),
 \qquad Q=V_\mu/L.
\]
The determinant square root and projected Taylor tensors give $\Omega^1$ and $K_X$; see 032, Lemma 4.1 and equations (4.5)--(4.8), pp.14--16. Proposition 4.6 and Corollary 4.7 vary the multiplicity-space data at compact places while keeping $V$ fixed. This operation does not change the signature of $V$.

\subsection{THETA3-B1: coherent rank-three definite neighbors}\label{B1}

\textbf{Statement.} There is a global Hermitian space $V^-$ with signature $(0,3)$ at $\mu$ and $(3,0)$ elsewhere, locally isometric to $V$ at every finite place. There is no everywhere positive definite $V^+$ with all those finite local isometry classes. For any finite nonsplit place $v_0$, there is an everywhere positive definite $V^{(v_0)}$ locally isometric to $V$ outside $\{\mu,v_0\}$, with the opposite determinant norm class at $v_0$.

\textbf{Proof.} Let $\epsilon_v$ be the local quadratic character of the determinant. At a real place it is $(-1)^{q_v}$, where $q_v$ is the negative index. A coherent collection has $\prod_v\epsilon_v=1$. Local dimension and determinant classify Hermitian spaces at finite places; together with the specified real signatures, the product condition gives global existence by Landherr's theorem. An exact source is B.H. Gross, \emph{Incoherent definite spaces and Shimura varieties}, 25 February 2018, \S3, Theorem 3.1, pp.6--7.

Changing the negative index at $\mu$ from one to three preserves $\epsilon_\mu$, so all finite data can be retained. Changing it from one to zero reverses $\epsilon_\mu$, contradicting the product condition if all other local data are retained. Reversing the determinant class at one nonsplit finite place reverses one more sign and restores the product condition. Landherr's theorem gives the asserted global spaces. \hfill$\square$

This is a statement about Hermitian isometry data and oscillator embeddings. Multiplying an odd-dimensional Hermitian form by a scalar can leave its abstract unitary group unchanged. Such a group identification does not identify the dual-pair embedding, splitting, or test vector.

\subsection{THETA3-B2: the scalar type distinguishes the definite neighbors}\label{B2}

\textbf{Statement.} Keep the active-place $\delta$ and multiplicity-space sign used in 032, and put $r=t^2$ on the double cover of the scalar circle of $\U(P_\mu)$. In the covariant Taylor model of 032, the degree-$k$ cohomological tensor has weight $t^{-3k}$. For the negative definite replacement in B1 the polynomial model
\[
 \bigoplus_{j\geq0}(\det\mathcal F^-)^{1/2}\otimes
                    \operatorname{Sym}^j\mathcal F^-
\]
has only weights $t^{3k+2j}$. The required weight is absent. For the positive definite replacement, its vacuum has weight $t^{-3k}$.

\textbf{Proof.} The scalar $r$ acts by $r$ on the $k$-dimensional summand $L\otimes P_\mu$ and by $r^{-1}$ on the $2k$-dimensional summand $Q^*\otimes P_\mu^*$. Hence $\det\mathcal F_P$ has weight $r^{-k}$, and its chosen square root has weight $t^{-k}$. The projected degree-$k$ tensor contributes $r^{-k}=t^{-2k}$, giving $t^{-3k}$.

Write $\delta_\mu=ia$, with $a\in\R^*$. The prescribed sign of $h_P$ is $-\operatorname{sgn}(a)$, as in 032, Lemma 4.1. When $h_V$ is negative definite, $a h_Vh_P$ is positive definite. The positive Lagrangian is then the nonconjugate summand of dimension $3k$, on which the scalar acts by $r$. Its half determinant contributes $t^{3k}$ and polynomial degree $j$ contributes $t^{2j}$. None gives $t^{-3k}$. When $h_V$ is positive definite, the positive Lagrangian is the conjugate summand, with scalar $r^{-1}$; its vacuum therefore has weight $t^{-3k}$. \hfill$\square$

The actual infinite Schr\"odinger coefficient vectors carry the dual action in 032, equation (4.13). Inverting both displayed spectra preserves both conclusions. A common honest character twist also preserves the mismatch. The same names for two splitting characters do not establish the same correction from the bare metaplectic action. This computation does not assert that an unspecified global theta packet is empty. Positive definite vacuum compatibility is necessary for this particular tensor comparison, not a sufficient identification of periods.

\subsection{THETA3-B3: an explicit local epsilon obstruction}\label{B3}

Fix a nonsplit quadratic local extension $E/F$, an additive character, $\delta$, a skew-Hermitian line $W$, a splitting character $\chi_W$, and a character $\alpha$ of the norm-one torus. If $H_+$ and $H_-$ are the two Hermitian lines of opposite determinant classes, then exactly one of the corresponding lifts of $\alpha$ to $W$ is nonzero.

Indeed, Borade--Franzel--Girsch--Yao--Yu--Zelingher, \emph{On tori periods of Weil representations of unitary groups}, Selecta Math. 31 (2025), article 49, Theorem 3.5, gives the condition
\[
 \epsilon_{E/F}(\chi_W^{-1}\alpha_E,\psi,\delta)
       =\epsilon(H)\epsilon_\delta(W),
 \qquad \alpha_E(x)=\alpha(x/\overline x).
\]
The left side is fixed and the two right sides are opposite signs, so exactly one equality holds. This proves the assertion as an application of that theorem. The source's \S\S3.2.1 and 4.2.1 state the splitting compatibilities before the seesaw identities. The example shows why local signs cannot be suppressed merely because a global parameter or a branching dimension agrees. It does not identify every Alp\"oge Bessel period with this particular torus period.

\subsection{THETA3-B4: normalization is indispensable for a numerical bound}\label{B4}

\textbf{Statement and proof.} Nonvanishing alone gives no positive lower bound independent of scalar normalization. If a nonzero theta input is replaced by $\varepsilon$ times itself, with $\varepsilon>0$, the resulting sesquilinear norm or simultaneous mixed pairing is multiplied by $\varepsilon^2$. It remains nonzero and tends to zero. Likewise, a nonzero constant determinant-line trivialization can be rescaled. Thus every numerical lower bound must specify the input norms and period conventions. \hfill$\square$

\subsection{THETA3-B5: a quantitative compact projection certificate}\label{B5}

\textbf{Statement.} Let $M$ be an $N$-dimensional unitary representation of a compact group $H$, with probability Haar measure. Let
\[
 P=\int_H\rho(h)\,dh
\]
have rank $r>0$. Some $r\times r$ principal minor of $P$ is at least $\binom Nr^{-1}$. In particular, when $r=1$, a normalized invariant functional has squared value at least $1/N$ on some vector of any orthonormal basis.

\textbf{Proof.} Haar invariance gives $P^2=P$; unitarity and invariance under inversion give $P^*=P$. Thus $P$ is the orthogonal projection onto $M^H$. Choose an $N\times r$ matrix $U$ whose columns form an orthonormal basis of that subspace. Then $P=UU^*$ and $U^*U=I_r$. Cauchy--Binet gives
\[
 1=\det(U^*U)=\sum_{|I|=r}|\det U_I|^2.
\]
There are $\binom Nr$ terms. For some $I$, its term is at least the reciprocal of this number, and that term is $\det P_{I,I}$. For $r=1$ these are exactly the squared coefficients of a normalized invariant vector, or equivalently of its normalized invariant functional. \hfill$\square$

For compact branching, take $M=\pi\otimes\sigma^\vee$. A multiplicity-one statement then gives a selected normalized period of squared size at least $1/(\dim\pi\dim\sigma)$. This does not set an unnormalized global automorphic integral equal to one.

\subsection{THETA3-B6: iteration requires an overlap estimate}\label{B6}

\textbf{Statement.} Positivity and nonzero eigenvalues of individual stages do not bound their composite. If instead
\[
 \|T_jw\|\geq b_j\|w\|,\qquad b_j>0,
\]
holds on the image of $T_{j-1}\cdots T_1$, then
\[
 \|T_m\cdots T_1v\|\geq\left(\prod_{j=1}^m b_j\right)\|v\|.
\]

\textbf{Proof.} In $\C^2$, let $P$ project onto $e_1$ and $Q_\theta$ onto $(\cos\theta,\sin\theta)$. Both have rank one and nonzero eigenvalue one, but $\|Q_\theta P\|=|\cos\theta|$, which can be arbitrarily small or zero. Under the stated additional hypothesis, apply the bound for $T_1$, then that for $T_2$ to $T_1v$, and continue. Each argument lies in the specified image, and multiplying the inequalities proves the result. \hfill$\square$

Thus multiplicity gained at each theta-lift stage does not, by itself, give a lower bound for a fixed composite period.

\subsection{THETA3-B7: an arithmetic Gram-matrix lower bound}\label{B7}

\textbf{Statement.} Let $K$ be CM with real subfield $K_0$ of degree $d$. Let $G\in M_r(K)$ be Hermitian and positive definite at every embedding $\sigma:K_0\hookrightarrow\R$. Suppose $DG$ has algebraic-integer entries for a positive integer $D$, and
\[
 \operatorname{tr}(G_\sigma)\leq T_\sigma,\qquad T_\sigma>0.
\]
For a chosen $\sigma_0$ and $r\geq2$,
\[
 \lambda_{\min}(G_{\sigma_0})\geq
 D^{-rd}\prod_{\sigma\ne\sigma_0}\left(\frac r{T_\sigma}\right)^r
       \left(\frac{r-1}{T_{\sigma_0}}\right)^{r-1}.
 \tag{B7}
\]
For $r=1$, omit the last factor. The right side also bounds
$a^*G_{\sigma_0}a/\|a\|_2^2$ from below for nonzero $a$.

\textbf{Proof.} The determinant $\Delta=\det G$ lies in $K_0$. The number $D^r\Delta$ is a totally positive algebraic integer, so its norm to $\Q$ is a positive integer and at least one. Consequently
\[
 \prod_\sigma\det G_\sigma\geq D^{-rd}.
\]
For $\sigma\ne\sigma_0$, the arithmetic--geometric mean inequality gives
$\det G_\sigma\leq(T_\sigma/r)^r$. Divide to obtain the corresponding lower bound for $\det G_{\sigma_0}$. The product of the other $r-1$ eigenvalues at $\sigma_0$ is at most $(T_{\sigma_0}/(r-1))^{r-1}$, so division gives (B7). When $r=1$ the determinant is the single eigenvalue. The final assertion follows from the spectral theorem. \hfill$\square$

The bounds at the other embeddings are essential: the totally positive unit $(3-2\sqrt2)^m$ tends to zero at one real embedding while its conjugate tends to infinity. On a finite definite adelic module, algebraic-valued vectors and rational mass weights can supply such a matrix once the field, denominator, and conjugate bounds are specified. For cohomological periods one must likewise identify an algebraic normalization and the relevant positivity at all embeddings; arbitrary complex theta coordinates do not have these properties automatically.

\subsection{THETA3-B8: a conditional two-place transfer inequality}\label{B8}

Suppose compatible normalized positive period formulas have been established:
\[
 I_{\mathrm{ball}}=C_{\mathrm{ball}}L\prod_{v\in S}Z_{\mathrm{ball},v},
 \qquad
 I_{\mathrm{def}}=C_{\mathrm{def}}L\prod_{v\in S}Z_{\mathrm{def},v},
\]
with the same nonzero $L$-factor and identical local factors outside
$\{\mu,v_0\}$. If
\[
 I_{\mathrm{def}}\geq b,\qquad
 \frac{C_{\mathrm{ball}}}{C_{\mathrm{def}}}\geq c>0,
 \qquad \frac{Z_{\mathrm{ball},v}}{Z_{\mathrm{def},v}}\geq a_v>0
 \quad(v=\mu,v_0),
\]
then
\[
 I_{\mathrm{ball}}\geq bc\,a_\mu a_{v_0}.
\]
The proof is division of the two identities followed by cancellation and multiplication of the displayed lower bounds. This proves an implication; it does not supply the two identities for the particular 032 and Alp\"oge inputs.

An applicable Rallis statement must be checked in its actual direction and range. Gan--Qiu--Takeda, arXiv:1207.4709v3, \S11.3, p.52, displays the doubling seesaw and its determinant-character twist. Section 11.4 assumes $d(n)<m\leq2d(n)$ before Proposition 11.1, including its convergent clause, and Theorem 11.4 repeats that range. For unitary groups $d(n)=n$, so it does not directly cover the naive $\U(1)\to\U(3)$ substitution. Dividing a local zeta integral by its standard local $L$-factor can still leave the normalizing denominator $d_v(s,\chi_V)$; one must not silently declare all such spherical factors to be one.

An applicable positive locator for the $\U(1)\to\U(3)$ direction is J. Haan, arXiv:1403.5061v7, \S3.3, Theorem 3.4, p.8. With Hermitian source dimension one, skew-Hermitian target dimension three, splitting $(\gamma,\gamma^3)$, $\gamma|_{\A_F^*}=\chi_{E/F}$, and the paper's Tamagawa and symplectic conventions, it states
\[
 \langle\theta(\overline f_1,\phi_1),\theta(\overline f_2,\phi_2)\rangle
 =\frac{L_E(3/2,\operatorname{BC}(\tau)\otimes\gamma^3)}
        {L_F(3,\chi_{E/F})}
   \prod_v Z_v^{\#}(1,f_{1,v},f_{2,v},\Phi_{1,v}).
\]
Here $Z_v^{\#}=L_v(3,\chi_{E_v/F_v})Z_v/L_{E_v}(3/2,\operatorname{BC}(\tau_v)\gamma_v^3)$ and the doubled section comes from $\rho_{1,3}(\phi_1\otimes\overline\phi_2)$. The source proves it by the convergent formula of J.-S. Li; it is not an application of the restricted GQT theorem just audited. Transport of the splitting, the conjugated source input, and the measures is still required. The formula controls an individual torus theta lift; it does not automatically give the norm of an arbitrary wedge of two such lifts. Haan's main Bessel theorem has additional special-pair hypotheses.

The recovered archive also retains a quotient-volume factor explicitly. In
\texttt{Development/\allowbreak{}AutomorphicArithmetic/\allowbreak{}ThetaCorrespondence/\allowbreak{}RallisInnerProduct/\allowbreak{}GraphRelationExtendsFromFockVectorsToKFinite.lean}, lines 791--802, the proof of \texttt{rallisInnerProductConvergent\_of} passes through the identity
\[
 \frac{\langle\Theta(\Phi_1,f_1),\Theta(\Phi_2,f_2)\rangle_B}
      {\operatorname{vol}([B])}
 =\operatorname{dblInt}_{A,\chi_B^{-1}}
       (\operatorname{ThetaMean}_B(\delta(\Phi_1,\Phi_2));f_1,f_2).
\]
Lines 776--789 explicitly convert the integral over $[B]$ to its volume times the normalized theta mean. The proof uses the conjugated source functions in its Fubini step; its lift convention is not interchangeable with \texttt{liftS}. Indeed, \texttt{GlobalTheory.RallisIP.thetaLift}, defined in \texttt{PublishedTheorems/\allowbreak{}TwistedTraceFormula/\allowbreak{}LowerLiftsZeroOfWittIdxEqZero.lean}, lines 22--26, is $h\mapsto\int_{[A]}\theta(g,h;\Phi)\overline{f(g)}\,dg$; lines 41--59 state the resulting range and volume-normalized zeta identity. Its input \texttt{Published.SiegelWeilUnitaryConvergent}, defined in \texttt{PublishedTheorems/\allowbreak{}SiegelWeil$_2$/\allowbreak{}SiegelWeilUnitaryConvergent.lean}, lines 18--35, requires
\[
 \operatorname{rank}A<\operatorname{rank}B\leq2\operatorname{rank}A,
 \qquad \operatorname{rank}A+\operatorname{wittIdx}B<\operatorname{rank}B.
\]
Thus this archive formula does include the anisotropic adjacent-rank alternatives $1\to2$ and $2\to3$ of B15.

A \emph{different} law, \texttt{SphericalVolumeLift.Convergent}, in \texttt{BlockTheory/\allowbreak{}TwistedDiscreteSeries/\allowbreak{}SphericalVolumeLift.lean}, lines 75--91, imposes $\operatorname{rank}B=\operatorname{rank}A+t$ and $\operatorname{rank}A<t$. Its range is $\operatorname{rank}B>2\operatorname{rank}A$; it is this larger-rank law that the recovered definite-positivity argument uses. Neither of these two ranges includes equal rank. There is also a separate \texttt{Published.SiegelWeilAnisotropicUnitary} statement in \texttt{PublishedTheorems/\allowbreak{}TwistedTraceFormula/\allowbreak{}SiegelWeilAnisotropicUnitary.lean}, lines 109--119: with anisotropic $B$ and positive ranks, its Eisenstein value is $(1\text{ if }\operatorname{rank}A<\operatorname{rank}B,\,2\text{ otherwise})$ times the normalized theta mean. This supplies an analytic identity at equal rank, not automatic nonvanishing of a given theta lift. The volume, character, and rank hypotheses must be kept separately for each formula.

\subsection{THETA3-B9: the canonical local Bessel integral equals one}\label{B9}

The recovered \texttt{Formalism/\allowbreak{}PeriodIntegrals/\allowbreak{}BesselPeriodGGP.lean} defines a real local nonvanishing proposition. Its ambient matrix \texttt{Hpk} has signature $(k+r,r+1)$; the compact subgroup \texttt{UWset} fixes the hyperbolic coordinates and a negative vector. Its integral has the form
\[
 \int_{\U(W_z)\times N}
 \langle\sigma(x_0)w,\sigma(hn)\sigma(x)w\rangle
       \overline{\psi_N(n)}\,\pi(h)^{-1}\,dh\,dn.
 \tag{B9a}
\]
It is operator-valued in the finite-dimensional $\pi$-space. The proposition asks for a discrete-series $\sigma$, a lowest-$K$-type unit vector, compact $\pi$, interlacing and half-integral differences, and asserts existence of $x,x_0$ making the integral nonzero. It does not say that arbitrary periods are zero or one. \texttt{GapBig} is in a separate numerical setting, not in the actual proposition's assumptions.

\textbf{Statement.} Put $k=2$, $r=0$, and $z=e_3$, and use probability Haar measures. Suppose a unitary representation of $\U(2,1)$ contains a unit vector $w$ with
\[
 \sigma(\operatorname{diag}(A,d))w=\det(A)d^{-2}w,
 \qquad A\in\U(2),\ d\in\U(1).
\]
Take $\pi(\operatorname{diag}(A,1))=\det A$ on $\C$. Then (B9a), at $x=x_0=1$, equals one.

\textbf{Proof.} The matrix equations defining \texttt{Nset} at $r=0$ force $N=\{1\}$. A unitary matrix fixing $e_3$ preserves its positive orthogonal complement, so \texttt{UWset} is precisely $\operatorname{diag}(\U(2),1)$. In the convention linear in the second inner-product variable,
\[
 \langle w,\sigma(\operatorname{diag}(A,1))w\rangle=\det A.
\]
Multiplication by $\pi(h)^{-1}$ cancels this factor. The integrand is identically one; compactness gives absolute convergence, and probability Haar gives the stated value. \hfill$\square$

This is the correct type for the right-regular coefficient of a canonical two-form. At a negative line $o$, the tangent space is $\Hom(L,Q)$. For $k=\operatorname{diag}(A,d)$, its derivative is $A/d$. If
\[
 f_\xi(g)=(g^*\xi)_o(e_1,e_2),
\]
then
\[
 f_\xi(gk)=\det(A)d^{-2}f_\xi(g).
\]
Left lattice invariance follows from that of $\xi$. On a compact quotient this coefficient has finite nonzero $L^2$ norm when $\xi\ne0$, and can be normalized. The displayed character is trivial on the center of $\U(2,1)$, so it is compatible with the projective special-unitary action. It is a coefficient convention; the action on the canonical-bundle fibre is dual.

The exact numerical choice in the archive is
\[
 \lambda=(1,0,-1),\quad v=(1),\quad b=(0,1),\quad
 u=(3/2,1/2),\quad\epsilon=(+,+,-).
\]
The definitions in \texttt{UpqRep.lean} give
\[
 \rho_c=(1/2,-1/2,0),\quad\rho_n=(1/2,1/2,-1),\quad
 \lambda+\rho_n-\rho_c=(1,1;-2).
\]
For compact $\U(2)$, $u-\rho_c=(1,1)$, so $\pi=\det$. The interlacing is
$3/2>1>1/2>0$; the Harish--Chandra integrality and half-integral differences follow directly from these tuples. The direct evaluation does not require the general discrete-series theorem. Invoking that theorem instead would additionally require the relevant \texttt{IsDSHpk} representation identification. No global period has been set equal to one.

The same direct calculation works for a pluricanonical coefficient of $K_X^q$, $q\geq1$. Its compact character is $(\det A)^qd^{-2q}$, and taking $\pi=\det^q$ again makes the local period one. The corresponding parameters are
\[
 \lambda=(q,q-1,1-2q),\quad v=q,\quad
 b=(1-q,2q-1),\quad u=(q+1/2,q-1/2).
\]
The same half-sum calculation gives lowest weight $(q,q;-2q)$, and the interlacing follows from
$q+1/2>q>q-1/2>q-1$. The node gap is $3q-2$. Thus powers can enlarge the numerical gap, whereas a determinant twist shifts both nodes equally and cannot do so. A section $\xi^q$ is pluricanonical; it is not the original ordinary cohomology class in $H^2(X)$. Passing a quantitative period bound back to that class needs an additional argument.

\subsection{THETA3-B10: a separate high-rank matching theorem}\label{B10}

The recovered \texttt{Development/\allowbreak{}AbelianVarietyTheory/\allowbreak{}BlockTheory/\allowbreak{}DivisorBase/\allowbreak{}WeilTensorModels.lean} defines one global matching theorem. It is distinct from the iterated Bessel family in B13. Its two spaces have dimensions $a,a'$, and its common target has dimension $a+a'$. The exact equality \texttt{IsConst}, line 491, is
\[
 I_V(\Phi)=c\,I_{V'}(\Phi')
\]
for matched data. \texttt{OneConst}, line 503, requires uniqueness and $c\ne0$. \texttt{StepAlpha} and \texttt{StepBeta}, lines 571 and 559, contain recursive constants involving partial Hecke $L$-values and local epsilon/Gauss factors. \texttt{LocalMatching} and \texttt{ThetaIntegralsAgree} occur at lines 581 and 585.

Its \texttt{Geo.IsGood}, lines 360--370, requires
\[
 a'<a,\quad a+a'\equiv0\pmod2,\quad
 2(p+q)+1<\left\lfloor\frac{a-a'}2\right\rfloor,\quad
 2(p+q)+1<a',
\]
and signatures
\[
 V:(a-p,p),\quad V':(q,a'-q),\quad
 W:(a+a'-p-q,p+q).
\]
The definitions \texttt{Spc.IsHerm} and \texttt{Spc.IsSkew}, lines 182--200, impose these signatures at every real place. The base also has CM/Galois and integral unramified-extension restrictions.

For $p=1,q=0$, the inequalities force $a'\geq4$ and the even difference $a-a'\geq8$; the smallest pair is $(12,4)$, with common target dimension sixteen. Thus assigning $a=3$ does not instantiate this theorem. Increasing the dimension of a one-active-place ball does not by itself repair its all-place signature or source-dimension requirements. This calculation audits this one matching theorem; it does not exclude the different definite-subgroup route or the abstract family theorem below.

\subsection{THETA3-B11: the exact restriction theta seesaw}\label{B11}

\textbf{Statement.} Let $V=V_H\perp L$ be Hermitian, with $L$ a line. Put $G=\U(V)$ and $H=\U(V_H)$, embedded by fixing $L$ pointwise. Let $W$ be a skew-Hermitian source and $T=\U(W)$. Assume $[H]$ and $[T]$ are compact, or assume absolute convergence of all sums and integrals below. Fix compatible oscillator splittings for the orthogonal tensor rule. In particular the splitting character for $V$ is the product of those for $V_H$ and $L$, and the source splittings agree.

For a finite sum of Schwartz tensors
$\Phi=\sum_a\Phi_{H,a}\otimes\Phi_{L,a}$ and a unitary automorphic character $\eta$ of $T$, set
\[
 F_\eta(g)=\int_{[T]}\theta_V(g,t;\Phi)\eta(t)^{-1}\,dt,
 \qquad
 B_e(F)=\int_{[H]}F(h)\overline{e(h)}\,dh.
\]
For a continuous automorphic test function $e$ on $[H]$, define
\[
 \Theta_{H\to T}(e;\Phi_H)(t)=
 \int_{[H]}\theta_{V_H}(h,t;\Phi_H)\overline{e(h)}\,dh.
\]
Then
\[
 B_e(F_\eta)=\sum_a\int_{[T]}
 \Theta_{H\to T}(e;\Phi_{H,a})(t)
 \theta_L(1,t;\Phi_{L,a})\eta(t)^{-1}\,dt.
 \tag{B11}
\]
All measures are the displayed measures, so no unspecified scalar occurs.

\textbf{Proof.} A polarization compatible with the orthogonal sum identifies the rational indexing lattice with the direct sum of the two indexing lattices. The tensor identity for the oscillator action identifies each summand for a pure tensor with a product. Absolute convergence permits separation of the sums, giving
\[
 \theta_V((h,1),t;\Phi_H\otimes\Phi_L)
       =\theta_{V_H}(h,t;\Phi_H)\theta_L(1,t;\Phi_L).
\]
Sum over $a$, multiply by $\overline{e(h)}\eta(t)^{-1}$, and integrate. Fubini applies by the convergence hypothesis. The line factor and $\eta$ are independent of $h$, and the inner $h$-integral is the displayed reverse theta lift. This proves (B11). In the compact case, uniform convergence on compact representative sets and boundedness of $e$ supply absolute integrability. \hfill$\square$

For a theta differential form, the statement applies to its polynomial Fock coefficient once that coefficient and its orthogonal tensor decomposition have been identified. This is a required identification of vectors, not merely of form degrees.

The actual archive declaration is
\texttt{GlobalTheory.T.thetaTensorRule\_orthogonalSum\_source}, in
\texttt{Proofs/\allowbreak{}DecoratedLatticeProofs/\allowbreak{}ThetaTensorRuleOrthogonalSumSourceStatement.lean}, with the proof in the matching file without \texttt{Statement}. The property \texttt{TensorSource}, in \texttt{BlockTheory/\allowbreak{}TwistedDiscreteSeries/\allowbreak{}ConjugateAdjointDagger.lean}, lines 27--53, explicitly imposes the splitting product and tensor action. The adjacent target tensor rule is used by \texttt{LastSeeSawThetaIdentity.lean}.

The precise global Bessel orientation is also recovered. In
\texttt{Definitions/\allowbreak{}TwistedInnerProducts/\allowbreak{}WittSmallGroupEmbeddingMemPointwiseStabilizerIso.lean}, lines 176--181, its integrand is
\[
 \phi(h)f(n\,\operatorname{emb}_A(h))\overline{\psi_{\rm Bes}(n)}.
\]
Its \texttt{liftS\_apply}, line 163, integrates $\theta(g,h)G(h)$ without conjugating $G$. Thus (B11) matches by taking $\phi=\overline e$ and $G=\eta^{-1}$. At Bessel rank zero the unipotent is trivial.

The source character $\eta$ remains visible in (B11). The identity therefore supplies a genuine global definite-theta/Bessel comparison, but does not on its own prove the cancellation claimed for the special iterated family. Nor does it equate the full ball norm with its restriction norm.

\subsection{THETA3-B12: every nonzero canonical form has a definite Bessel coefficient}\label{B12}

\textbf{Statement.} Let $X$ be the compact congruence ball surface of B0 and $0\ne\xi\in H^0(X,K_X)$. There is an $E$-rational negative line $L$ such that $H=\U(L^\perp)$ is totally definite of rank two and a finite-level automorphic test function $e$ with $B_e(f_\xi)\ne0$. This evaluates a canonical-bundle coefficient, not the cohomological pullback of $H^2(X)$ to a point.

More precisely, use probability Haar on $[H]$, a compact open finite level $K_H$ fixing $f_\xi$, and the finite-dimensional space $R$ with this level and the same one-dimensional infinity type. Put $N_H=\dim R$, let $m_0>0$ be the mass of the identity double coset, and let $b=|f_\xi(1)|$. For any orthonormal basis $(e_j)$ of $R$,
\[
 \sum_{j=1}^{N_H}|B_{e_j}(f_\xi)|^2
      =\|f_\xi|_{[H]}\|_2^2\geq m_0b^2,
 \qquad
 \max_j|B_{e_j}(f_\xi)|^2\geq\frac{m_0b^2}{N_H}.
 \tag{B12}
\]
In particular, $\xi\ne0$, $\int_X\xi\wedge\overline\xi>0$, and existence of some such nonzero Bessel coefficient are equivalent.

\textbf{Proof.} First make the adelic realization precise. Fix the $\mathcal O_E$-lattice defining the principal congruence quotient. An element of $\U(V)(F)$ preserving this lattice has determinant in $\mathcal O_E^\times$ with $\alpha\overline\alpha=1$. Every complex conjugate of $\alpha$ has absolute value one, so $\alpha$ is a root of unity. There are finitely many such roots in $E$. Choose a principal congruence level at which no nonidentity root is congruent to one. Then
\[
 \Gamma'=\U(V)(F)\cap K_f\subset\SU(V)(F).
\]
Make the level deeper so that $\Gamma'\subset\Gamma$, where $X=\Gamma\backslash D$; this inclusion has finite index. Thus $X'=\Gamma'\backslash D$ is both a finite principal cover of $X$ and a connected component of the adelic $\U(V)$ quotient. Pull $\xi$ to $X'$ and extend it by zero on the other components. This gives the automorphic canonical-bundle coefficient $f_\xi$ used below. Its central archimedean character is trivial, since $\det(tI_2)t^{-2}=1$. If the covering degree is $d$, then
\[
 \int_{X'}p^*\xi\wedge\overline{p^*\xi}
       =d\int_X\xi\wedge\overline\xi,
\]
so nonvanishing and positivity are unchanged. The masses and finite level in (B12) refer to this chosen realization. We retain the notation $\xi$ for its pullback.

The nonzero locus of $\xi$ is a nonempty open subset of the ball. Rational vectors are dense at the active embedding: the image of $E$ in $\C$ contains $\Q+\Q\alpha$ for a nonreal $\alpha$, hence is dense. In coordinates, $V(E)$ is therefore dense in $V_\mu$. Choose $z\in V(E)$ whose negative line lies in this open locus. At every other real place the form on $V$ is positive definite, so $h_\tau(z,z)>0$. Thus $L^\perp$ has signature $(2,0)$ everywhere and $H$ is totally definite.

Base the canonical coefficient at $L$, choosing a nonzero tangent two-frame there. The calculation of B9 gives the right compact character $\det(A)d^{-2}$, and on $H_\infty$, which fixes $L$ pointwise, it is $\det A$ at the active place. The other compact factors have the prescribed canonical coefficient type, trivial for the trivial coefficient system. Its absolute value is therefore constant along each right $H_\infty$ orbit. By construction $f_\xi(1)\ne0$. The adelic realization just constructed gives an open compact $K_H$ fixing the coefficient.

The set $H(F)\backslash H(\A_F)/(H_\infty K_H)$ is finite. On each double coset a function of the specified scalar infinity type is determined by one complex number, subject to its finite stabilizer condition. Thus $R$ is finite-dimensional. Let $m_i>0$ be the masses of these cosets. The identity coset contributes exactly $m_0b^2$ to the norm square, and all other terms are nonnegative. Consistent Haar choices give masses proportional to reciprocal finite stabilizer orders, normalized by $\sum_i m_i=1$. At neat level the stabilizers are trivial and the masses are equal.

Parseval in $R$ proves the equality in (B12); one term is at least their average. The orthonormal basis may be refined within the automorphic spectral summands of the compact quotient, so these are periods against definite automorphic test vectors. This assertion does not require their values or their Hilbert normalizations to be algebraic.

Conversely, a nonzero Bessel coefficient forces $f_\xi$, and hence $\xi$, to be nonzero. In local holomorphic coordinates, if $\xi=a\,dz_1\wedge dz_2$, then
\[
 \xi\wedge\overline\xi
   =4|a|^2\,dx_1\,dy_1\,dx_2\,dy_2.
\]
Its integral in the complex orientation is positive exactly when $\xi\ne0$. This proves the equivalence. \hfill$\square$

There is also a certified point-selection version. If a unit-frame coefficient has size $b_0>0$ at a point, has Lipschitz constant $L_\xi$ nearby, and the Hermitian form has a known negative margin, choose a rational point within the margin-preserving radius and within $b_0/(2L_\xi)$ when $L_\xi>0$. Then $b\geq b_0/2$ and the lower bound is $m_0b_0^2/(4N_H)$. THETA3-Q3 gives an explicit denominator from the coordinate matrix, margin, and radius. The resulting level, mass, and dimension still require bounds or computation.

For two restrictions in the same $R$, Parseval gives the mixed identity
\[
 \int_{[H]}f\overline{f'}=
 \sum_j B_{e_j}(f)\overline{B_{e_j}(f')}.
\]
THETA3-A5 passes from finitely many restrictions to the full ball inner product by a positive inverse frame operator. A naive quotient disintegration does not do this: it unfolds to $H(F)\backslash G(\A)$, rather than $G(F)\backslash G(\A)$. The rational or algebraic nature of the frame correction is a further issue.

Finally, a central lift of the Shimura cocharacter at $L$ may factor through $\U(L^\perp)\times\U(L)$ rather than $H$ alone. The proof uses the finite definite adelic module and canonical-fibre evaluation. It does not assert an unspecified morphism of Shimura data $H\to G$.

\subsection{THETA3-B13: the recovered character-independent functional}\label{B13}

\paragraph{The exact family.}
The source \texttt{Definitions/\allowbreak{}LatticeFrames/\allowbreak{}ValLastX.lean}, under
\texttt{Development/\allowbreak{}AbelianVarietyTheory}, contains the relevant definitions. The following locators separate data from conclusions:
\begin{itemize}
\item \texttt{Comm}, lines 28--57, fixes the indefinite spaces $V_j$, discrete-series models, definite spaces $X_j$, Bessel/Jacobi decompositions, measures, final $T$, and $X_m=T\perp L_z$.
\item \texttt{Tup}, lines 79--113, supplies varying Weil representations, finite translation combinations, row functions, and constants. \texttt{RealEq}, lines 182--184, requires their real Weil actions to agree.
\item \texttt{perF}, lines 119--120, is a finite sum of genuine Bessel periods of finite right translates. \texttt{per}, lines 122--123, inserts the row function and automorphic realization. \texttt{perT}, lines 137--138, is the final definite restriction period.
\item \texttt{up}, lines 166--167, iterates two definite theta operators with the intervening $S$ operation; \texttt{dn}, lines 169--170, composes the reverse automorphic lift maps.
\end{itemize}
The normalizing products, lines 172--174, are
\[
 c_{\rm Top}(n)=
 \left(\prod_{i=0}^{n-1}c_V(2i)c_K(2i)c_K(2i+1)\right)c_P,
 \qquad
 c_{\rm Bot}=\prod_{j=0}^{m-1}c_I(j).
\]
\texttt{LawsX}, lines 212--221, includes common-data validity, tuple validity, finite support, common real action, the double-step and last-step integral laws, definite kernel reduction, final period factorization, and row covariance.

The actual source witness is
\texttt{GlobalTheory.SeeSaw.mainX\_of\_lawsX}, in
\texttt{Proofs/\allowbreak{}DecoratedLatticeProofs/\allowbreak{}SeesawTransportRelationsOfFamilyLaws.lean}, lines 25--37. Its inputs are \texttt{CuspRapidDecay}, \texttt{HarishChandraModules}, and \texttt{LawsX}; its conclusion is \texttt{MainX}. The matching \texttt{Statement.lean}, line 12, exposes the wrapper
\texttt{GlobalTheory.T.seesawTransportRelations\_of\_familyLaws}.
These declarations and the displayed proof body were read, but their complete transitive dependencies were not audited or compiled.

\paragraph{Mathematical content.}
Fix $m=2n+2$. For an allowed tuple, input column, and terminal automorphic vector, indexed collectively by $s$, put
\begin{align*}
 A_s&=c_{{\rm Top},s}(n+1)
      \operatorname{tail}_s(\operatorname{up}_s(n+1)\psi_s,F_s),\\
 B_s(y,v)&=c_{{\rm Bot},s}
      \operatorname{per}_s\bigl(\psi_s,
       \operatorname{dn}_s(2n+1)\operatorname{ioT}_s(F_s);y,v\bigr).
\end{align*}
\texttt{MainX}, lines 244--251, says that every finite scalar relation
$\sum_s\alpha_s A_s=0$ implies
$\sum_s\alpha_s B_s(y,v)=0$ for every fixed $y,v$.
Consequently there is a single scalar $a(y,v)$ such that
\[
 B_s(y,v)=a(y,v)A_s
 \quad\text{for every allowed tuple, column, and terminal vector.}
 \tag{B13}
\]
The scalar is unique if some $A_s\ne0$. This is the precise normalized independence supplied by the recovered theorem. It is not independence under arbitrary changes of characters and Schwartz data.

\textbf{Proof of the consequence.} If all $A_s$ vanish, one-term relations force all $B_s(y,v)=0$, and $a=0$ works. Otherwise choose $s_0$ with $A_{s_0}\ne0$ and put $a=B_{s_0}(y,v)/A_{s_0}$. For each $s$, the two-term relation
$A_{s_0}A_s-A_sA_{s_0}=0$ gives
$A_{s_0}B_s-A_sB_{s_0}=0$. Division proves (B13), and also uniqueness. This is a complete elementary re-proof of the logical passage from relation transfer to the common scalar. It does not re-prove the analytic inputs of \texttt{MainX}. \hfill$\square$

\paragraph{The exact instantiation and arithmetic gates.}
The abstract family is more flexible than the particular tower used in the construction. In \texttt{Definitions/\allowbreak{}LatticeFrames/\allowbreak{}JointSmooth.lean}, lines 56--62, \texttt{DefSp.Is} requires a compact quotient and a Cartan involution equal to $+1$ or $-1$ at each real place, hence definiteness, without requiring the same sign everywhere. Its \texttt{BesAt}, lines 107--115, consists of a line orthogonal decomposition and a hyperbolic decomposition. Its rank-zero local geometry therefore includes the active $\U(2,1)/\U(2)$ pair of B9.

By contrast, the concrete \texttt{PiLocalization$_2$.Is} in
\texttt{Definitions/\allowbreak{}LatticeFrames/\allowbreak{}ChainParSucc.lean}, lines 56--60, requires the same signature \texttt{sigYD D.R j} at every embedding. In \texttt{ExistsDatOfTower.lean}, lines 151--171, the Bessel step is built with a globally negative line, of signature $(0,1)$ everywhere. Hence that constructed tower does not immediately equal the one-active-place 032 datum. The theorem's abstract shape does not itself supply a witness satisfying all these hypotheses, nor the discrete-series span and gap requirements used to construct the tower.

The gap obstruction can be stated numerically. In \texttt{BlockTheory/\allowbreak{}TwistedDiscreteSeries/\allowbreak{}ThetaSourceAction.lean}, lines 209--211, \texttt{NodeGap} requires a natural number $N\geq\operatorname{rank}X'+1$ such that
\[
 t+2N^3+1\leq b_{i+1}-b_i
\]
for every adjacent pair and every infinite place. Consequently its right-hand gap is at least three. The canonical choice $b=(0,1)$ in B9 cannot satisfy this condition when identified with the two nodes. This requirement is used to construct the tower's spanning data, for example in \texttt{StepVanOfStepX.lean}, lines 71--84; it is not an additional assumption of the direct local integral evaluated in B9, nor of the elementary consequence (B13) once \texttt{LawsX} has already been supplied.

The bridge to the arithmetic columns is explicitly named. In
\texttt{Definitions/\allowbreak{}LatticeFrames/\allowbreak{}RatioOKOfRootOfUnity.lean}, the property
\texttt{FamAll}, lines 60--62, combines \texttt{LawsX} with
\texttt{LinkZero}, \texttt{LinkTail}, \texttt{LinkDown},
\texttt{LinkOnto}, \texttt{LinkComm}, and \texttt{RatioOK}.
The last condition compares $c_{\rm Top}/c_{\rm Bot}$ across tuples.
\texttt{UnitRatio}, in \texttt{Definitions/\allowbreak{}LatticeFrames/\allowbreak{}UnitRatio.lean}, lines 22--23, means
\[
 y\ne0,\quad \exists n>0,\ r\in\Q^*:
 (x/y)^n=r,
 \qquad v_p(r)\ne0\Longrightarrow p\in S.
\]
Thus the source controls radical-of-rational ratios with specified prime support; it does not automatically identify their absolute values with one. To derive a quantitative period bound, one must instantiate these links for the cohomological vectors, retain the normalizers, and bound the relevant common scalar or normalized Gram matrix. B9 and B12 give explicit local and definite-subgroup nonvanishing inputs, while THETA3-A5 supplies a separate exact finite-frame comparison of the global mixed pairing.

\subsection{THETA3-B14: large weights, tensor contractions, and the return to cohomology}\label{B14}

The large-weight device has an exact scope. For $q\geq1$, the canonical-power coefficient has compact type $(\det A)^qd^{-2q}$ and the parameters displayed after B9. Its local Bessel integral is one and its node gap is $3q-2$. Thus, for a fixed numerical threshold $M$, the condition $3q-2\geq M$ is attained by $q\geq\lceil(M+2)/3\rceil$. This does not alone construct the remaining tower data, and changes the input from a two-form to a pluricanonical section.

\paragraph{A cohomological distinction.}
For $q>1$, the coherent multiplication map
\[
 \operatorname{Sym}^q H^0(X,K_X)\longrightarrow H^0(X,K_X^q)
\]
is not the exterior-product map followed by diagonal pullback. If
$0\ne\xi\in H^{2,0}(X)$, its external product lies in
$H^{2q,0}(X^q)$, whereas diagonal pullback lies in
$H^{2q,0}(X)=0$, since $2q>\dim_\C X=2$.
The coherent section $\xi^q$ is nonzero: in a local canonical frame it is the $q$th power of its coefficient, tensor the $q$th power of the frame. These are different operations.

There is no nonzero global holomorphic bundle map $K_X^q\to K_X$ on a compact ball surface. Indeed, $K_X$ is ample, and such a map would be a section of $K_X^{1-q}$. A nonzero section would give an effective divisor $D$ with
\[
 D\cdot K_X=(1-q)K_X^2<0,
\]
contrary to ampleness; the nowhere-zero case would make $K_X^{1-q}$ trivial, also contradicting $K_X^2>0$. This rules out a global fibrewise holomorphic contraction of that specified kind. It does not rule out other correspondences between auxiliary varieties or maps requiring extra data.

Division by $\xi^{q-1}$ is an inverse on the particular image
$\xi^{q-1}H^0(X,K_X)\subset H^0(X,K_X^q)$, because cancellation gives a holomorphic result on that image. On an arbitrary section it can have poles along the zero divisor of $\xi$. Thus using this inverse requires proof that the high-weight input belongs to this image. A nonzero vector produced somewhere in the full high-weight space does not supply that proof.

\paragraph{The precise tensor contraction and its nonvanishing.}
Let $V=H^2(X,\Q)$ and
\[
 Q_X(u,v)=\int_Xu\smile v:\ V\otimes V\longrightarrow\Q(-2).
\]
Put $A=Q_X(\xi,\overline\xi)>0$. Every contraction of two factors in the pure tensor $\xi^{\otimes q}$ by $Q_X$ vanishes, because $Q_X(\xi,\xi)=0$ by Hodge type. By contrast the contraction
\[
 C_q:V^{\otimes q}\otimes V^{\otimes(q-1)}
       \longrightarrow V(-2(q-1)),
\]
\[
 C_q(u_1\otimes\cdots\otimes u_q\otimes v_2\otimes\cdots\otimes v_q)
       =u_1\prod_{i=2}^q Q_X(u_i,v_i)
\]
satisfies
\[
 C_q(\xi^{\otimes q}\otimes\overline\xi^{\otimes(q-1)})
           =A^{q-1}\xi\ne0.
 \tag{B14a}
\]
There is no new nonvanishing hypothesis in this calculation beyond $\xi\ne0$.

The map $C_q$ is induced by algebraic diagonals and push-forward. On $X^{2q-1}$, pair the $i$th and $(q+i-1)$st factors, for $2\leq i\leq q$, by the pullback of the diagonal class of $X\times X$, and push to the first factor. For an external product, the projection formula gives one factor $\int_Xu_i\smile v_i$ from each pair. All factors have even degree, so no permutation sign occurs. The resulting Tate twist is exactly $-2(q-1)$. This proves both the formula and its algebraic-correspondence realization.

This construction needs an actual external-product cohomological input, including the conjugate factors. A coherent section of $K_X^q$ alone is not such an input. Nor does the fact that $C_q$ is algebraic turn a non-Hodge vector such as $\xi^{\otimes q}\otimes\overline\xi^{\otimes(q-1)}$ into an algebraic cycle. In a proposed cycle construction, the ambient Hodge types and the required projected component must first be identified; after that, algebraic contraction preserves algebraicity and (B14a) supplies the appropriate nonzero test value.

A parallel representation-theoretic check applies to a specified rational Weil component. Suppose a polarizing algebraic torus acts on
$W_\C=\ell_\chi\oplus\ell_{\chi^{-1}}$, where $\chi$ has infinite order, and its invariant pairing pairs the two lines. Any torus-equivariant map from $W^{\otimes q}$ to $W$ sends the pure weight $\chi^q$ to zero for $q>1$, since neither $\chi$ nor $\chi^{-1}$ equals $\chi^q$. Pairwise polarization contractions are equivariant and therefore have this obstruction. With $q-1$ opposite-character factors the total weight is $\chi$, and the contraction is nonzero whenever their pairings are nonzero. This is a statement about that explicitly specified torus and equivariant contractions; it is not an impossibility theorem for every algebraic correspondence. It pinpoints which conjugate component a tensor-power argument must retain.

\paragraph{Analytic return with an explicit norm loss.}
Fix a Hermitian metric on $K_X$ and a positive measure on $X$. Write
\[
 A_q=\int_X|\xi|^{2q},\qquad M=\|\xi\|_\infty>0.
\]
Pointwise $|\xi|^{2q}\leq M^{2q-2}|\xi|^2$, hence
\[
 A_1\geq \frac{A_q}{M^{2q-2}}.
 \tag{B14b}
\]
Therefore a certified lower bound for the global norm of the particular pure power $\xi^q$, together with a certified upper bound for $M$, gives a lower bound for the original norm. A local Bessel integral normalized to one is not by itself that global bound.

There is also a nonzero analytic contraction. Use the metric to contract
$\xi^q\otimes\overline{\xi}^{\otimes(q-1)}$ to the smooth canonical section
$|\xi|^{2q-2}\xi$, and take its orthogonal projection to $H^0(X,K_X)$. Its inner product with $\xi$ is $A_q>0$, so the projection is nonzero. This projection depends on the metric and the original form and is not thereby an algebraic correspondence. Equations (B14a) and (B14b) give two valid return mechanisms with their required input data; neither identifies an arbitrary high-weight output with the original Weil component automatically.

\subsection{THETA3-B15: a small-rank local alternative and its exact global gates}\label{B15}

There is an alternative to the large-weight construction: the canonical discrete series itself has a nonzero local theta correspondence with a compact rank-two group. We give the characters and parameters explicitly. This removes a local occurrence obstruction; it does not identify a particular global theta vector with the iterated lift of B13.

\paragraph{The literal one-active-place datum forces the rank-zero step.}
Suppose that $V_0$ in \texttt{Comm} is the rank-three space $V$ of B0 and that $m=2n+2$. Then $n=0$, hence $m=2$. Indeed, \texttt{Comm.bes}, at \texttt{ValLastX.lean}, line 46, requires a \texttt{BesAt} datum of rank $n$. By \texttt{JointSmooth.lean}, lines 107--115, this contains a global hyperbolic decomposition of $V$ with $n$ hyperbolic planes. A nonzero isotropic vector over $E$ remains a nonzero isotropic vector at every completion. Since $V$ is positive definite at another real place, it has no such vector. Its global Witt index is zero, and therefore $n=0$.

At $m=2$ all the intermediate conditions indexed by $j+2(\sigma+1)+2=m$ are empty. The rank-zero \texttt{JacAt} identifies $V_1$ with the definite space $X_1$; write this middle space as $S$. Thus the remaining geometric pattern is
\[
 V=A\perp L,\qquad S=X_1\text{ totally definite},\qquad
 X_2=T\perp L_z\text{ totally definite}.
 \tag{B15a}
\]
Here $A$ is the definite rank-two complement in B12. The abstract \texttt{Comm} does not impose the concrete ranks \texttt{rankY} on its spaces $V_j$. Its field \texttt{par}, line 33, may vary with the real place. The numerical \texttt{Ranks} conditions do allow $m=2$: the choices $k=6,t=13,\operatorname{ord}=1$ satisfy the inequalities in \texttt{Formalism/\allowbreak{}LocalRingCardinality/\allowbreak{}Params.lean}, lines 6--22. This numerical observation does not construct the other fields of \texttt{D7} or a full \texttt{Comm.Is} witness.

The last seesaw imposes a precise line relation, not an arbitrary real signature change. In \texttt{PlaceLaws.lean}, lines 165--166, \texttt{WeilOK.tau} is
\[
 \operatorname{rank}L_z=1,\qquad
 \tau^*G_L\tau=-N_{E/F}(\xi_0)G_{L_z}.
 \tag{B15b}
\]
The norm is positive at every real place. Hence $L_z$ has the opposite sign to $L$: positive at $\mu$, negative elsewhere. This has a coherent global realization: set $L_z=-N_{E/F}(\xi_0)^{-1}L$ and $\tau=1$. Taking $T=L_z$ gives a rank-one top and a definite rank-two $X_2=T\perp L_z$, with positive sign at $\mu$ and negative sign elsewhere. These are actual global forms; no unverified collection of finite determinant classes is being prescribed. Choosing them does not yet prove the needed finite local theta occurrences.

\paragraph{Convention and local theorem.}
Use Ichino's convention in \emph{Theta lifting for tempered representations of real unitary groups}, arXiv:2008.06174v2, \S4.1, p.10. The input space $W$ is skew-Hermitian of dimension $n$ and signature $(p,q)$, the target space is Hermitian of dimension $r+s=m$, and
\begin{equation}
\begin{gathered}
 \chi_V(z)=(z/|z|)^{m_0},\qquad
 \chi_W(z)=(z/|z|)^{n_0},\\
 m_0\equiv m\pmod2,\qquad n_0\equiv n\pmod2,
 \qquad \psi(x)=e^{-2\pi i x}.
\end{gathered}
 \tag{B15c}
\end{equation}
The local lift is characterized by
$\Hom_{\U(W)\times\U(V)}(\omega,\pi\boxtimes\theta(\pi))\ne0$.
Let $\pi_\mu$ have Harish--Chandra parameter $(1,0;-1)$ on $\U(2,1)$, and let $\pi_\tau=1$ on $\U(3)$ for $\tau\ne\mu$. For the pair of dimensions $(n,m)=(3,2)$, take $m_0=-2,n_0=-1$. Then
\[
 \theta_{2,0}(\pi_\mu)=\det,\qquad
 \theta_{2,0}(\pi_\tau)=\det\quad(\tau\ne\mu).
 \tag{B15d}
\]
Here the character $\det$ of compact $\U(2)$ has Harish--Chandra parameter $(3/2,1/2)$.
The representation $\pi_\mu$ is the holomorphic discrete series with lowest compact type $(1,1;-2)$, namely $(A,d)\mapsto\det(A)d^{-2}$. This is the canonical right-regular coefficient convention fixed in B9: $\lambda+\rho_n-\rho_c=(1,1;-2)$. Replacing the right-action convention by its dual would require dualizing this parameter and the subsequent pairings as well.

For the next pair, let $\sigma=\det^{-1}$ on $\U(2)$, with Harish--Chandra parameter $(-1/2,-3/2)$. The following local lifts are nonzero and trivial:
\[
\begin{array}{c|c|c|c|c}
\text{place}&(n,m)&(m_0,n_0)&\text{target signature}&\theta(\sigma)\\ \hline
\mu&(2,1)&(-3,-2)&(1,0)&1\\
\tau\ne\mu&(2,1)&(-1,2)&(0,1)&1.
\end{array}
 \tag{B15e}
\]
Thus there are compatible middle representations for a two-step composition at every real place: the first oscillator has a quotient $\pi_\tau\boxtimes\det$, and the second has a quotient $\det^{-1}\boxtimes1$.

\textbf{Proof.}
In all three calculations the rank drops by one. After subtracting $m_0/2$ from the input Harish--Chandra parameter, the nonvanishing criterion has a single zero chain, of length $k=1$. For (B15d) the shifted parameters are
\[
 (2,1;0)\quad\text{and}\quad(2,1,0),
\]
respectively. In Atobe's notation both have $(k_\lambda,r_\lambda,s_\lambda)=(1,2,0)$. Atobe, \emph{On the non-vanishing of theta liftings of tempered representations of $\U(p,q)$}, Adv. Math. 363 (2020), Definition 1.6 and Theorem 1.7(2), pp.4--5, gives nonvanishing at the base point $(2,0)$. The theorem for discrete series is unconditional, as stated immediately after it. Ichino, \S6.2, explicitly identifies Atobe's additive-character convention with (B15c).

The parameter rule for adjacent ranks deletes the zero, transfers positive entries in the first block and negative entries in the second block to the positive target block, and then adds $n_0/2$. This is the adjacent-rank case of Ichino's Theorem 4.1(ii), pp.11--12, already covered by Paul's adjacent-rank correspondence cited there. In (B15d) it gives $(2,1)-(1/2,1/2)=(3/2,1/2)$. For the first row of (B15e), the shifted parameter is $(1,0)$, giving base point $(1,0)$ and output $1-1=0$. For the second row it is $(0,-1)$, giving base point $(0,1)$ and output $-1+1=0$. The same nonvanishing theorem applies in both cases. Every parity condition in (B15c) holds. Compact rank-one Harish--Chandra parameter zero is the trivial representation. This proves the claims. \hfill$\square$

The abstract groups in this calculation must be identified with the Hermitian/skew-Hermitian spaces of 032 and the archive, including the sign of the imaginary scalar used to pass between the forms. Equations (B15c)--(B15e) specify the convention rather than silently making that identification. Likewise, the archive's \texttt{liftS} integrates the kernel against the input function bilinearly. A quotient $\pi\boxtimes\det$ therefore pairs with an automorphic input of type $\det^{-1}$ to produce $\pi$. This is why the middle input in (B15e) is the dual of the B9 compact test representation. The two pairs are allowed different splitting characters; \texttt{WeilOK.last}, lines 160--164, separately names the second pair's character $\chi'_S$.

\paragraph{The archimedean characters can be made global.}
Let $E/F$ be CM, let $\epsilon$ be its quadratic idele-class character, and fix integers $a_\tau$ all congruent to $d$ modulo two. There is a unitary Hecke character $\chi$ of $E$ satisfying
\[
 \chi|_{\A_F^\times}=\epsilon^d,
 \qquad \chi_\tau(z)=(z/|z|)^{a_\tau}.
 \tag{B15f}
\]
Finite ramification is allowed; prescribed finite local values are not part of this assertion.

\textbf{Proof.}
The natural map $C_F\to C_E$ is a closed embedding. Indeed it is injective, and the idele norm of $C_E$ restricts to the square of the norm on $C_F$; the norm-one idele-class group of $F$ is compact, so the embedding is proper. The character-extension theorem for closed subgroups of locally compact abelian groups supplies a unitary extension $\chi_0$ of $\epsilon^d$. This is precisely the argument in 032, Lemma 4.4, using Folland, \emph{A Course in Abstract Harmonic Analysis}, second edition, \S4.3, Corollary 4.42. Since $\chi_0$ is trivial on each positive real scalar, its archimedean characters are $(z/|z|)^{a^0_\tau}$, with $a^0_\tau\equiv d\pmod2$.

Put $T_1=\operatorname{Res}^{1}_{E/F}\mathbf G_m$. Its real group $T_1(F\otimes\R)$ embeds as a compact subgroup of $T_1(F)\backslash T_1(\A_F)$ by placing one at the finite components: the intersection with the diagonal rational group is trivial. Character extension on this locally compact abelian quotient gives an automorphic character $\rho$ with archimedean weights $(a_\tau-a^0_\tau)/2$. Then
\[
 \chi(z)=\chi_0(z)\rho(z/\overline z)
\]
is trivial on $E^\times$, has the same restriction to $\A_F^\times$ as $\chi_0$, and has the prescribed weights. Its finite components have open kernels on compact unit groups, so it has finite conductor. \hfill$\square$

This proves that the integer choices in (B15d)--(B15e) themselves have no global parity obstruction. It does not force the finite components of those characters to be those in a previously fixed theta construction.

\paragraph{What changing a splitting transports.}
For a fixed dual pair, changing the splitting multiplies the oscillator action by determinant characters of the same two groups. Locally, Atobe's Proposition 3.3 and Proposition 3.5, pp.20--22, give
\begin{equation}
\begin{aligned}
 \omega_{m_0+2a}|_{\U(p,q)}
       &\simeq\omega_{m_0}|_{\U(p,q)}\otimes\det^a,\\
 \theta^{m_0+2a}(\pi)
       &\simeq\theta^{m_0}(\pi\otimes\det^{-a}),
\end{aligned}
 \tag{B15g}
\end{equation}
when the other splitting character is kept fixed. The latter formula includes zero lifts. A simultaneous change of the other character tensors the output by its corresponding determinant character.

The archive exposes its orientation in
\texttt{Formalism/\allowbreak{}LocalToGlobal/\allowbreak{}IndefiniteReal/\allowbreak{}LocalNewformFunction.lean}, line 51:
\[
 \texttt{splitFacO}(\chi_S,\chi_T;x,y)=\chi_T(x)^{-1}\chi_S(y).
 \tag{B15h}
\]
If $\kappa_S(z)=\rho_S(z/\bar z)$ and $\kappa_T(z)=\rho_T(z/\bar z)$ are character ratios, its split-place scalar changes by
$\rho_T(\det g)^{-1}\rho_S(\det h)$.
The determinants here follow from \texttt{TiedWeil.lean}, lines 71--78, and the actual \texttt{IsWeilDatumO.om\_split} condition at lines 72--75 of \texttt{LocalNewformFunction.lean}. Thus an inverse in one convention cannot simply be removed when comparing it with (B15g).

All these determinant factors restrict trivially to the appropriate special unitary group. The canonical $\SU(V)$ geometry can consequently survive a determinant twist of its $\U(V)$ realization. However, the underlying source space is unchanged: a mixed-sign rank-two $P$ in A7 remains that same $P$, and a local group $\U(1,1)$ does not become $\U(2)$ by a character twist. Compensating the target twist as in (B15g) preserves the original local correspondence. To replace the source by a different definite inner form while keeping a fixed target vector, one must establish occurrence in the new dual pair. The original A7 realization does not alone provide that occurrence at every finite place, or an automorphic representation on the new source.

\paragraph{A precise global nonvanishing test for the alternative.}
There is a usable criterion with no unspecified real-place correction. In the conventions of Gan--Qiu--Takeda, arXiv:1207.4709v3, Theorem 11.7, pp.56--57, let $U_n$ be the source of dimension $n=1$ or $2$, and let $V$ be an anisotropic target of dimension $m=n+1$. For an irreducible cuspidal automorphic source representation $\Sigma$, its lift to this particular $V$ is nonzero if and only if
\begin{equation}
\begin{gathered}
 \theta_{V_v}(\Sigma_v)\ne0\quad\text{for every place }v,
 \\ L(1,\Sigma\times\chi_V)\ne0
       \quad\text{with a holomorphic value}.
\end{gathered}
 \tag{B15i}
\end{equation}
Here $L$ is the source's completed standard doubling $L$-function, with the character as named in that paper. The hypotheses are $d(n)=n<m=n+1\leq2n$, target Witt index $r=0$, and hence no lower step in its global Witt tower. The equality $m=d(n)+1$ is exactly the final case in Theorem 11.7(ii); Proposition 11.6 says that it removes the otherwise possible real-sign replacement. No temperedness hypothesis is needed for this criterion.

With Tamagawa measures and the paper's normalized doubling factors, Theorem 11.4(ii), p.54, specializes to
\begin{equation}
\begin{gathered}
 \langle\theta(\phi_1,f_1),\theta(\phi_2,f_2)\rangle
 \\
 =2L(1,\Sigma\times\chi_V)
       Z^*(1/2,\phi_1\otimes\overline{\phi_2},f_1,f_2),
 \\
 Z_v^*(s)=\frac{Z_v(s)}{L(s+1/2,\Sigma_v\times\chi_{V,v})}.
\end{gathered}
 \tag{B15j}
\end{equation}
The global product includes the unramified denominator factors of that normalization; they must not be replaced by one. The archive's bilinear input and Ichino's character names must be converted before substituting a representation in this formula. For comparison, Ichino's \S8.2, p.37, uses an incomplete $L$-function twisted by $\chi_V^{-1}$ and, for dimensions $2\to3$ and tempered source, divides it by $L^S(2,1)L^S(3,\epsilon)$ while retaining the local matrix-coefficient integrals at the excluded places. These are different displayed normalizations of the standard method, not an asserted equality of the named characters.

Thus both arrows in a proposed rank-one to rank-two to rank-three global construction have a specific value-at-one criterion, together with finite local occurrence. This is not a central-value criterion. It still must be shown that the resulting rank-three representation contains the prescribed canonical theta vector, rather than merely some canonical vector. For a fixed rank-three representation, the reverse first-occurrence question is the pole side of the Rallis theory; an upward criterion for an independently chosen $\Sigma$ does not establish that reverse occurrence automatically.

The definite side of the small choice in (B15a) has a different rank pattern: $A_2\to S_2\to X_2$, with $\operatorname{rank}X_2=2$. Its two arrows have equal rank. The adjacent-rank criterion (B15i) does not cover them; their standard Rallis parameter would be the central point, and their nonvanishing requires its own argument. In particular, the strict range of the archive's \texttt{SphericalVolumeLift.Convergent} just recorded in B8 is not satisfied. No claim that every definite-side lift is nonzero is part of the local calculation.

\paragraph{The universal identity and the arithmetic equation left to prove.}
The last-step analytic identity is already a universal source theorem, rather than a large-gap instance. In \texttt{Proofs/\allowbreak{}DecoratedLatticeProofs/\allowbreak{}DoubleLiftPeriodRankZeroFubiniCompactQuotient.lean}, line 136, the source proves \texttt{lastLaw : \ensuremath{\forall} D, LastLaw D}; in \texttt{LastSeeSawThetaIdentity.lean}, line 22, it proves \texttt{seeLast : \ensuremath{\forall} D, SeeLast D}. The proof bodies use compact quotient integration and the exact tensor identity with (B15b). With the requisite Weil data and
\[
 \mu_{V_1}=(\operatorname{emb}_0)_*\mu_{X_1},
 \tag{B15k}
\]
which is \texttt{Comm.VolLast} in \texttt{HzUnit.lean}, lines 23--24, the theorem \texttt{GlobalTheory.LastRed.main}, in \texttt{LastTwoStepsDoubleLiftPeriodEqIntegral.lean}, line 293, gives
\begin{equation}
\begin{gathered}
 B_A\bigl(\varphi,
   \Theta_{V\leftarrow S}(\Theta_{S\leftarrow T}(F;\Phi_2);\Phi_1)\bigr)
 \\
  =\int_{[T]} F(t)\,
       \operatorname{biLift}_{A\to S\to X_2}
       (\varphi;C_{\Phi_1,\Phi_2})(\iota(t,1))\,dt.
\end{gathered}
 \tag{B15l}
\end{equation}
This is the untranslated specialization; the source additionally gives every finite/right translate explicitly. The contracted kernel is its actual \texttt{lastBi}, defined in \texttt{CombAtSingleZero.lean}, lines 29--34: on the decomposition $X_2=T\perp L_z$ it is
\[
 C_{\Phi_1,\Phi_2}(u;v_T,v_z)
       =\Phi_2(v_T)\Phi_1(u,\operatorname{ofAdj}_{L,S}(\tau v_z)).
\]
All characters and measures remain in this identity. There is no arbitrary comparison scalar at this stage. The required global input is the concrete double theta lift appearing on its left.

The additional equation needed to identify these analytic kernels with the arithmetic columns is explicit. In \texttt{Definitions/\allowbreak{}LatticeFrames/\allowbreak{}OnFramesRdApply.lean}, lines 108--113, the condition \texttt{ClsK.Core} is
\begin{equation}
\begin{gathered}
 \operatorname{volLev}(X_j,U_{j,\alpha})\,
 \theta_{K,j}(h_{j,\alpha,a},h_{j+1,\gamma,a'};\Phi_{j,P})
 \\
 =c_{K,j}\,\mathcal A_{j,\alpha,\gamma,a,a'}
       \bigl(\operatorname{mvZ}(1,s_{j,\tau}I)(\operatorname{castP}P)\bigr).
\end{gathered}
 \tag{B15m}
\end{equation}
Here $\mathcal A$ is the specified \texttt{abelian\ensuremath{_2}} functional, not an arbitrary linear functional defined to make the equality true. Its class sums carry the reciprocal stabilizer orders in lines 115--121. The terminal condition \texttt{ClsT.Is.tail}, lines 147--150, similarly includes the top-level covolume and stabilizer masses. Their conjunction with the real branching and level conditions is \texttt{ClsPinD}, lines 152--153, from which the source derives the corresponding \texttt{KerX} and \texttt{PerX} laws. The universal Fubini identity (B15l) does not prove (B15m), supply those arithmetic functionals, or establish that the given 032 vector lies in the indicated double-lift image. These are now concrete representation and kernel/covolume equations, rather than an unspecified appeal to global matching.
